\section{Introduction}
\label{sec:intro}

Text-to-image diffusion transformers such as FLUX.1~\citep{flux2024} are trained on images of about one megapixel. Sampled directly at $4096\times4096$, they repeat objects, tile textures at the training period, and lose fine detail~\citep{he2024scalecrafter,du2024demofusion,zhang2025diffusion4k}. Two mechanisms inside attention change at once. RoPE~\citep{su2024roformer} sees query-key offsets up to four times longer than any seen in training, and each softmax normalizes over sixteen times as many tokens.

Training-free methods address the first mechanism by remapping RoPE frequencies. Position interpolation~\citep{chen2023pi} compresses all of them, NTK-aware scaling and YaRN~\citep{peng2024yarn} treat bands differently, LongRoPE~\citep{ding2024longrope} searches per-band factors, and DyPE~\citep{issachar2025dype} makes the remapping depend on the diffusion time. They address the second mechanism with a global logit temperature, either YaRN's $0.1\ln s+1$ or a temperature that grows with the logarithm of the token count~\citep{jin2023training,nakanishi2025ssmax}. SEGA~\citep{sega2026} adds a third ingredient: a separate q/k gain for every RoPE dimension, set at each step from the Fourier spectrum of the noisy latent. Because the attention logit is a Fourier series in the query-key offset (Section~\ref{sec:background}), such a gain acts as an equalizer on the positional kernel of attention.

SEGA's gain law is not derived from a stated objective, and our analysis (Section~\ref{sec:sega}) shows what it computes. Its allocation z-scores the log band energies across RoPE dimensions, which makes it invariant to the slope of a power-law spectrum: a steep and a shallow spectrum receive the same tilt, which differs only through one scalar flatness gate. In mechanism tests of the paper equations on forward-noised real latents, a single static tilt explains 71\% of the variance of its modulation and image content at most 9\%. The released implementation also differs from the paper equations in input features, frequency mapping, amplitude, and clipping (Table~\ref{tab:sega-diff}). These observations leave open whether a per-band gain helps because it adapts to content or because it applies a fixed bias toward local bands.

We propose a per-band gain law with a stated reason for each term. Sampling at $s$ times the training resolution resembles viewing a picture from $s$ times closer. The human visual system meets this change whenever we approach an image, and psychophysics has measured how it keeps perceived contrast and focus stable. We translate five of these findings into the DiT setting (Section~\ref{sec:method}). Efficient coding under noise~\citep{atick1990towards,atick1992retina,vanhateren1992theory} gives a Wiener visibility gate per band, computed against the noise floor of the flow, which is known exactly at every step. Contrast constancy~\citep{georgeson1975contrast} and blur adaptation~\citep{webster2002neural} give a correction proportional to the log amplitude deficit relative to a natural spectrum, which we take to be the spectrum the model itself produces at its native resolution. Octave-band channels~\citep{campbell1968application,field1987relations} set the band structure, and divisive normalization~\citep{heeger1992normalization,carandini2012normalization} motivates the log-token-count temperature that anchors all gains. The result, Constancy RoPE, is a closed loop: the gain at step $t$ depends on the deficit of the clean-image prediction from step $t-1$, and it vanishes when that prediction already has the native spectral profile.

Our contributions are:
\begin{itemize}
  \item An analysis of SEGA's gain law covering the paper equations and the released code. We prove that its allocation is invariant to the slope of a power-law spectrum and to any power transform of the spectrum, show that the released code maps RoPE dimensions to spectral bins by index rank rather than wavelength, and report mechanism tests on forward-noised real latents (Section~\ref{sec:sega}).
  \item Constancy RoPE, a training-free per-band q/k gain derived from efficient coding, contrast constancy, blur adaptation, octave channels, and divisive normalization. It has one gain parameter, a fixed point at the native spectrum, a per-band visibility gate, and a sign that follows from the derivation (Section~\ref{sec:method}).
  \item An evaluation protocol for FLUX.1-dev at up to $4096^2$ with metrics that see 4K detail (patch FID and KID on native-resolution crops, crop-level no-reference quality, spectral slope error, and a tiling index), prompt-level bootstrap confidence intervals, and falsification tests for both SEGA's claim and ours (Section~\ref{sec:setup}).
  \item \todo{Results contribution: one sentence stating what the head-to-head and kill tests found, including negative outcomes.}
\end{itemize}
